\section{Additional Implementation Details}
\label{appendix-details}


For all methods, we use the Adam~\cite{kingma2014adam} optimizer with $\beta_1 =$ 0.9 and $\beta_1 = $ 0.999, and a batch size of 128 images. We resize all images to 224x224 and normalize them to [0,1]. We train CIFAR-100 and DomainNet for 20 epochs, and ImageNet-R for 50 epochs (chosen to ensure models converge fully for each task). As discussed in the main text, we use the same prompting lengths and locations for L2P~\cite{wang2022learning} and DualPrompt~\cite{wang2022dualprompt} as recommended by the more recent DualPrompt paper. Specifically, for Dualprompt, we use a length 5 prompt in layers 1-2 (referred to as \emph{general} prompts) and length 20 prompts in layers 3-5 (referred to as \emph{task-expert}). For L2P, we use a prompt pool of size 20, total prompt length of size 20, and choose 5 prompts from the pool to use during inference.

As done in DualPrompt~\cite{wang2022dualprompt}, we tuned all additional hyperparameters using 20\% of the training data as validation data. This resulted in using a learning rate of $1e^{-3}$ for all prompting methods (as opposed to $5e^{-3}$ as reported in DualPrompt), and a learning rate of $1e^{-4}$ for all methods which fully fine-tune the model. We searched for learning rates in the values of $\{1e^{-6},5e^{-6},1e^{-5},5e^{-5},1e^{-4},5e^{-4},1e^{-3},5e^{-3},1e^{-2}\}$. We also found that cosine-decaying learning rate outperforms a constant learning rate (which was used in the original DualPrompt implementation). We conjecture that the reduced and decaying learning rate explain the performance boost we obtained on the 10-task ImageNet-R benchmark using our implementations.

For our method, we use a prompt length of 8 and 100 prompt components (and prompt at the same locations as DualPrompt), which were chosen with a hyperparameter sweep on validation data to have the best trade-off between performance and parameter efficiency. We searched for prompt lengths in the range of [4,40] and prompt components in the range of [5,500]. We use $\lambda = $ 0.1 to weight the orthogonality regularization loss, chosen from sweeping across decade values from $1e^{-6}$ up to $1e^2$. As shown in \AppA{}, we see that increasing the prompt length has little effect on our method, whereas increasing the prompt component size has strong returns all the way up to 200 components.

Finally, when implementing classification loss for fine-tuning, L2P, DualPrompt, and CODA-P, we re-use a technique from the official GitHub repo for the DualPrompt and L2P papers~\cite{wang2022learning,wang2022dualprompt} and replace the predictions from past-task logits with negative infinity when training a new task. This results in a softmax prediction of ``0" for these past task classes and prevents gradients from flowing to the linear heads of past task classes. While not discussed in these papers, this technique is \emph{crucial} for performance of these methods, as we confirmed during reproduction. Essentially, the linear layer is highly biased towards new tasks in class-incremental learning in the absence of rehearsal, so this technique prevents the linear head from learning a bias towards new classes over past classes. We note that this bias is a well-known issue~\cite{wu2019large,Ahn2021ssil}.

\section{Additional Results}
\label{appendix-metrics-and-results}

\begin{table*}[t]
\caption{\textbf{Results (\%) on 5-task ImageNet-R (40 classes per task)}. $A_N$ gives the accuracy averaged over tasks, $F_N$ gives the average forgetting, and $N_{param}$ gives the \% of trainable parameters and final parameters w.r.t. the base ViT pre-trained model. We report the mean and standard deviation over 5 trials.}
\vspace{-2mm}
\centering
\label{tab_2:imnet-r_5}
\begin{tabular}{c||c|c|c} 
\hline
\rule{0pt}{10pt} Method  & $A_N$ ($\uparrow$) & $F_N$ ($\downarrow$) & \thead{$N_{param}$ ($\downarrow$) \\ Train/Final}\\
\hline
Upper-Bound  & $77.13$ & - & $100/100$  \\ 
\hline
ER (5000)    & $71.72 \pm 0.71$ & $13.70 \pm 0.26$ & $100/100$  \\ 
FT           & $18.74 \pm 0.44$ & $41.49 \pm 0.52$ & $100/100$  \\
FT++         & $60.42 \pm 0.87$ & $14.66 \pm 0.24$ & $100/100$  \\
LwF.MC       & $74.56 \pm 0.59$ & $4.98 \pm 0.37$ & $100/100$  \\ 
L2P++        & $70.83 \pm 0.58$ & $3.36 \pm 0.18$ & $ 0.7/100.7$  \\
Deep L2P++   & $73.93 \pm 0.37$ & $2.69 \pm 0.10$ & $ 9.6/109.6$  \\ 
DualPrompt   & $73.05 \pm 0.50$ & $\bm{2.64 \pm 0.17}$ & $ 0.5/100.5$  \\
\hdashline
\textbf{CODA-P-S}  & $75.19 \pm 0.47$ & $2.65 \pm 0.15$ & $ 0.7/100.7$  \\
\textbf{CODA-P}    & $\bm{76.51 \pm 0.38}$ & $2.99 \pm 0.19$ & $ 4.6/104.6$  \\

\hline
\end{tabular}

\vspace{9mm}

\caption{\textbf{Results (\%) on 10-task ImageNet-R (20 classes per task)}. $A_N$ gives the accuracy averaged over tasks, $F_N$ gives the average forgetting, and $N_{param}$ gives the \% of trainable parameters and final parameters w.r.t. the base ViT pre-trained model. We report the mean and standard deviation over 5 trials.}
\vspace{-2mm}
\centering
\label{tab_2:imnet-r_10}
\begin{tabular}{c||c|c|c} 
\hline
\rule{0pt}{10pt} Method & $A_N$ ($\uparrow$) & $F_N$ ($\downarrow$) & \thead{$N_{param}$ ($\downarrow$) \\ Train/Final}\\
\hline
Upper-Bound  & $77.13$ & - & $100/100$  \\  
\hline
ER (5000)    & $64.43 \pm 1.16$ & $10.30 \pm 0.05$ & $100/100$  \\ 
FT           & $10.12 \pm 0.51$ & $25.69 \pm 0.23$ & $100/100$  \\ 
FT++         & $48.93 \pm 1.15$ & $9.81 \pm 0.31$ & $100/100$  \\
LwF.MC       & $66.73 \pm 1.25$ & $3.52 \pm 0.39$ & $100/100$ \\
L2P++        & $69.29 \pm 0.73$ & $2.03 \pm 0.19$ & $ 0.7/100.7$  \\
Deep L2P++   & $71.66 \pm 0.64$ & $1.78 \pm 0.16$ & $ 9.6/109.6$  \\  
DualPrompt   & $71.32 \pm 0.62$ & $1.71 \pm 0.24$ & $ 0.8/100.8$  \\
\hdashline
\textbf{CODA-P-S}  & $73.93 \pm 0.49$ & $\bm{1.60 \pm 0.20}$ & $ 0.7/100.7$  \\
\textbf{CODA-P}    & $\bm{75.45 \pm 0.56}$ & $1.64 \pm 0.10$ & $ 4.6/104.6$  \\

\hline
\end{tabular}

\vspace{9mm}
\caption{\textbf{Results (\%) on 20-task ImageNet-R (10 classes per task)}. $A_N$ gives the accuracy averaged over tasks, $F_N$ gives the average forgetting, and $N_{param}$ gives the \% of trainable parameters and final parameters w.r.t. the base ViT pre-trained model. We report the mean and standard deviation over 5 trials.}
\vspace{-2mm}
\centering
\label{tab_2:imnet-r_20}
\begin{tabular}{c||c|c|c} 
\hline
\rule{0pt}{10pt} Method & $A_N$ ($\uparrow$) & $F_N$ ($\downarrow$) & \thead{$N_{param}$ ($\downarrow$) \\ Train/Final}\\
\hline
Upper-Bound  & $77.13$ & - & $100/100$  \\ 
\hline
ER (5000)    & $52.43 \pm 0.87$ & $7.70 \pm 0.13$ & $100/100$  \\ 
FT           & $4.75 \pm 0.40$ & $16.34 \pm 0.19$ & $100/100$  \\ 
FT++         & $35.98 \pm 1.38$ & $6.63 \pm 0.11$ & $100/100$  \\
LwF.MC       & $54.05 \pm 2.66$ & $2.86 \pm 0.26$ & $100/100$  \\
L2P++        & $65.89 \pm 1.30$ & $1.24 \pm 0.14$ & $ 0.7/100.7$  \\
Deep L2P++   & $68.42 \pm 1.20$ & $1.12 \pm 0.13$ & $ 9.6/109.6$  \\ 
DualPrompt   & $67.87 \pm 1.39$ & $1.07 \pm 0.14$ & $ 1.3/101.3$  \\
\hdashline
\textbf{CODA-P-S}  & $70.53 \pm 1.24$ & $1.00 \pm 0.15$ & $ 0.7/100.7$  \\
\textbf{CODA-P}    & $\bm{72.37 \pm 1.19}$ & $\bm{0.96 \pm 0.15}$ & $ 4.6/104.6$  \\

\hline
\end{tabular}
\end{table*}








\begin{table*}[t]
\caption{\textbf{Results (\%) on 10-task CIFAR-100 (10 classes per task)}. $A_N$ gives the accuracy averaged over tasks, $F_N$ gives the average forgetting, and $N_{param}$ gives the \% of trainable parameters and final parameters w.r.t. the base ViT pre-trained model. We report the mean and standard deviation over 5 trials.}
\vspace{-2mm}
\centering
\label{tab_2:cifar}
\begin{tabular}{c||c|c|c} 
\hline
\rule{0pt}{10pt} Method & $A_N$ ($\uparrow$) & $F_N$ ($\downarrow$) & \thead{$N_{param}$ ($\downarrow$) \\ Train/Final}\\
\hline
Upper-Bound  & $89.30$ & - & $100/100$  \\  
\hline
ER (5000)    & $76.20 \pm 1.04$ & $8.50 \pm 0.37$ & $100/100$  \\ 
FT           & $9.92 \pm 0.27$ & $29.21 \pm 0.18$ & $100/100$  \\  
FT++         & $49.91 \pm 0.42$ & $12.30 \pm 0.23$ & $100/100$  \\  
LwF.MC       & $64.83 \pm 1.03$ & $5.27 \pm 0.39$ & $100/100$  \\ 
L2P++        & $82.50 \pm 1.10$ & $1.75 \pm 0.42$ & $ 0.7/100.7$  \\
Deep L2P++   & $84.30 \pm 1.03$ & $\bm{1.53 \pm 0.40}$ & $ 9.5/109.5$  \\  
DualPrompt   & $83.05 \pm 1.16$ & $1.72 \pm 0.40$ & $ 0.7/100.7$  \\
\hdashline
\textbf{CODA-P-S}  & $84.59 \pm 0.87$ & $1.76 \pm 0.28$ & $ 0.6/100.6$  \\
\textbf{CODA-P}    & $\bm{86.25 \pm 0.74}$ & $1.67 \pm 0.26$ & $ 4.6/104.6$  \\

\hline
\end{tabular}

\vspace{9mm}

\caption{\textbf{Results (\%) on 5-task DomainNet (69 classes per task)}. $A_N$ gives the accuracy averaged over tasks, $F_N$ gives the average forgetting, and $N_{param}$ gives the \% of trainable parameters and final parameters w.r.t. the base ViT pre-trained model. We report the mean and standard deviation over 3 trials.}
\vspace{-2mm}
\centering
\label{tab_2:domainnet}
\begin{tabular}{c||c|c|c} 
\hline
\rule{0pt}{10pt} Method  & $A_N$ ($\uparrow$) & $F_N$ ($\downarrow$) & \thead{$N_{param}$ ($\downarrow$) \\ Train/Final}\\
\hline
Upper-Bound  & $79.65$ & - & $100/100$  \\ 
\hline
ER (5000)    & $58.32 \pm 0.47$ & $26.25 \pm 0.24$ & $100/100$  \\
FT           & $18.00 \pm 0.26$ & $43.55 \pm 0.27$ & $100/100$  \\ 
FT++         & $39.28 \pm 0.21$ & $44.39 \pm 0.31$ & $100/100$ \\
LwF.MC       & $\bm{74.78 \pm 0.43}$ & $5.01 \pm 0.14$ & $100/100$  \\
L2P++        & $69.58 \pm 0.39$ & $2.25 \pm 0.08$ & $ 0.9/100.9$  \\
Deep L2P++   & $70.54 \pm 0.51$ & $2.05 \pm 0.07$ & $ 9.7/109.7$  \\  
DualPrompt   & $70.73 \pm 0.49$ & $\bm{2.03 \pm 0.22}$ & $ 0.6/100.6$  \\
\hdashline
\textbf{CODA-P-S}    & $71.80 \pm 0.57$ & $2.54 \pm 0.10$ & $ 0.6/100.6$  \\
\textbf{CODA-P}      & $73.24 \pm 0.59$ & $3.46 \pm 0.09$ & $ 4.8/104.8$  \\

\hline
\end{tabular}

\vspace{9mm}

\caption{\textbf{Results (\%) on ImageNet-R \emph{with covariate domain shifts}}. Results are included for 5 tasks (40 classes per task). We simulate domain shifts by randomly removing 50\% of the dataset's domains (e.g., clipart, paintings, and cartoon) for the training data of each task (see SM for more details). $A_N$ gives the accuracy averaged over tasks, $F_N$ gives the average forgetting, and $N_{param}$ gives the \% of trainable parameters and final parameters w.r.t. the base ViT pre-trained model. We report the mean and standard deviation over 5 trials.}
\vspace{-2mm}
\centering
\label{tab_2:imnet-r_domainshift}
\begin{tabular}{c||c|c|c} 
\hline
\rule{0pt}{10pt} Method & $A_N$ ($\uparrow$) & $F_N$ ($\downarrow$) & \thead{$N_{param}$ ($\downarrow$) \\ Train/Final}\\
\hline
Upper-Bound  & $77.13$ & - & $100/100$  \\  
\hline
ER (5000)    & $67.39 \pm 0.37$ & $11.94 \pm 0.17$ & $100/100$  \\  
FT           & $17.93 \pm 0.27$ & $37.49 \pm 0.28$ & $100/100$  \\
FT++         & $54.51 \pm 0.68$ & $14.41 \pm 0.33$ & $100/100$  \\
LwF.MC       & $64.02 \pm 1.55$ & $7.05 \pm 0.27$ & $100/100$  \\
L2P++        & $65.08 \pm 0.29$ & $2.79 \pm 0.32$ & $ 0.7/100.7$  \\
Deep L2P++   & $65.74 \pm 0.12$ & $2.48 \pm 0.30$ & $ 9.6/109.6$ \\  
DualPrompt   & $66.98 \pm 0.08$ & $\bm{2.21 \pm 0.28}$ & $ 0.8/100.8$  \\
\hdashline
\textbf{CODA-P-S}  & $69.73 \pm 0.18$ & $2.35 \pm 0.19$ & $ 0.7/100.7$  \\ 
\textbf{CODA-P}    & $\bm{71.35 \pm 0.08}$ & $2.56 \pm 0.26$ & $ 4.6/104.6$  \\

\hline
\end{tabular}
\end{table*}










We report extended results, including standard deviations and additional parameters trained, for all benchmarks in Tables \ref{tab_2:imnet-r_5} (5-task ImageNet-R~\cite{hendrycks2021many,wang2022dualprompt}), \ref{tab_2:imnet-r_10} (10-task ImageNet-R), \ref{tab_2:imnet-r_20} (20-task Imagenet-R), \ref{tab_2:cifar} (10-task CIFAR-100~\cite{krizhevsky2009learning}), \ref{tab_2:domainnet} (5-task DomainNet~\cite{peng2019moment}), and \ref{tab_2:imnet-r_5} (Dual-Shift ImageNet-R).

We evaluate methods using (1) final average accuracy $A_{N}$, or the final model's test accuracy averaged over all $N$ tasks, and (2) average forgetting~\cite{chaudhry2018efficient,Lopez-Paz:2017,mai2022online} $F_{N}$, or the drop in task performance averaged over $N$ tasks. The reader is referred to Appendix C of Wang \emph{et al.}~\cite{wang2022dualprompt} for the formal metric definitions.
We emphasize that $A_{N}$ is the more important metric and encompasses both method plasticity \emph{and} forgetting, 
whereas $F_{N}$ provides additional context subject to the model's plasticity (i.e., a lower $F_{N}$ value \emph{and} a lower $A_{N}$ value would indicate that the model's lower forgetting results from its reduced adaptivity to new tasks, which is an undesirable trait).

For each result, we calculate the mean and standard deviation over separate runs. Each run contains different shuffles of the class order; specifically, we shuffle the classes using a random seed that is set for each ``trial run" - and form the tasks using this class shuffle. \emph{Importantly, the class order and all randomized seeds, including model initialization, are consistent between different methods in the same ``trial run".}

We additionally report the number of parameters \emph{trained} (i.e., unlocked during training a task) as well as the \emph{total} number of parameters in the final model. These are reported in \% of the backbone model for easy comparison. Importantly, we design CODA-P-S to have fewer parameters than DualPrompt in the 10-task ImageNet-R setting (our main experiment setting). As we change the number of tasks in ImageNet-R, the number of total parameters for DualPrompt changes because the pool size is set as equal to the number of total tasks by definition (unlike ours, which is set as a hyper-parameter, allowing us to increase or decrease the number of trainable parameters to accommodate the underlying complexity of the task.)

\section{ImageNet-R Dual-Shift Benchmark}
\label{appendix-benchmark}

Our motivation for the challenging \emph{dual-shift} ImageNet-R~\cite{hendrycks2021many,wang2022dualprompt} benchmark is to show robustness to two different types of continual distribution shifts: semantic \emph{and} covariate. Specifically, there are 15 image types in the ImageNet-R dataset: \emph{`art', `cartoon', `deviantart', `embroidery', `graffiti', `graphic', `misc', `origami', `painting', `sculpture', `sketch', `sticker', `tattoo', `toy', `videogame'}. We divide the dataset into 5 tasks of 40 classes each, and within each task we randomly remove 8 of the domain types from the training data. The task becomes more challenging because we now have image type domain shifts injected into the continual learning task sequence, in addition to the already-present class-incremental shifts.
